\documentclass[11pt,a4paper]{article}

\usepackage[T1]{fontenc}
\usepackage[utf8]{inputenc}
\usepackage[brazil]{babel}
\usepackage{lmodern}
\usepackage[margin=2.5cm]{geometry}
\usepackage{booktabs}
\usepackage{array}
\usepackage{xcolor}
\usepackage{listings}
\usepackage{pgfplots}
\usepackage{hyperref}

\pgfplotsset{compat=1.18}
\hypersetup{colorlinks=true, urlcolor=blue, linkcolor=black}

\definecolor{fundo}{gray}{0.96}
\definecolor{borda}{gray}{0.80}

\lstset{
  basicstyle=\ttfamily\footnotesize,
  backgroundcolor=\color{fundo},
  frame=single,
  rulecolor=\color{borda},
  columns=fullflexible,
  keepspaces=true,
  showstringspaces=false,
  xleftmargin=0.5em,
  xrightmargin=0.5em,
  aboveskip=0.8em,
  belowskip=0.8em
}

\setlength{\parindent}{0pt}
\setlength{\parskip}{0.6em}
\setlength{\emergencystretch}{3em}

\newcommand{\cod}[1]{\texttt{#1}}
\newcommand{\separador}{\begin{center}\rule{0.45\linewidth}{0.4pt}\end{center}}


\title{Trabalho Prático de Agentes: Contract Net Protocol}
\author{Felipe Escobar Koizimi\\[0.3em]\small Sistemas Multiagentes -- Prof. Jomi F. Hübner}
\date{}

\begin{document}

\maketitle

Há $n$ agentes que precisam contratar um serviço (os \emph{initiators}) e $m$ agentes que oferecem esse serviço (os \emph{participants}). Cada initiator faz $i$ contratações ao mesmo tempo, com $n$ entre 2 e 199, $m$ entre 2 e 49 e $i$ entre 1 e 9. Todos os participants oferecem um único serviço, mas cada um define o seu preço conforme a sua estratégia.

Para resolver o problema, o Contract Net Protocol (CNP) foi implementado em Jason 3.3, com os agentes escritos em AgentSpeak e o ambiente em Java. O ambiente não participa das negociações: ele apenas registra os eventos e encerra a execução quando todas terminam. Também foram desenvolvidos scripts para executar os testes e os experimentos.

\begin{center}
\begin{tabular}{l>{\raggedright\arraybackslash}p{5.2cm}}
\toprule
Seção & Conteúdo \\
\midrule
1 & O protocolo \\
2 & Modelagem \\
3 & Método dos experimentos \\
4--6 & Experimentos 1 a 3: variação de n, m e i \\
7 & Experimento 4: estresse \\
8 & Avaliação da abordagem de agentes \\
9 & Avaliação das ferramentas \\
10 & Critério de comparação entre implementações \\
11 & Conclusões \\
\bottomrule
\end{tabular}
\end{center}

O código, os scripts e os dados usados neste relatório estão em \url{https://github.com/escobar-felipe/tp-cnp-jason}.

\separador

\section*{1. O protocolo}

Uma negociação do CNP tem cinco passos:

\begin{lstlisting}
initiator                         participants
    |--- cfp(Cid, servico_a) --------->|   (para todos)
    |<-- propose(Cid, Preco) ----------|   ou refuse(Cid)
    |   espera todos responderem e escolhe o menor preco
    |--- accept_proposal(Cid) -------->|   (vencedor)
    |--- reject_proposal(Cid) -------->|   (os outros)
    |<-- inform_done(Cid) -------------|   (vencedor, depois de executar)
\end{lstlisting}

O initiator anuncia o serviço para todos os participants (\cod{cfp}, de \emph{call for proposals}). Cada um responde com um preço ou com uma recusa. Quando todos respondem, ou quando o prazo de 2000\,ms termina, o initiator escolhe o menor preço, informa o vencedor e rejeita as demais propostas. O vencedor executa o serviço e informa a conclusão.

Quando os $m$ participants respondem, uma negociação troca $m$ mensagens de \cod{cfp}, $m$ de \cod{propose}, 1 de \cod{accept\_proposal}, $m$ $-$ 1 de \cod{reject\_proposal} e 1 de \cod{inform\_done}, ou seja, 3$m$ + 1 mensagens. Uma rodada completa troca $n$ $\times$ $i$ $\times$ (3$m$ + 1). Essa fórmula foi usada para verificar a contagem feita durante as execuções.

\separador

\section*{2. Modelagem}

\subsection*{2.1 Agentes e o sistema}

O sistema é descrito no \cod{cnp.mas2j}:

\begin{lstlisting}
MAS cnp {
    environment: cnp.CnpEnv(2, 3, 2)
    agents:
        initiator   [beliefs="tarefas(3)"] #2;
        participant                        #2;
    aslSourcePath: "src/agt";
}
\end{lstlisting}

O \cod{\#2} cria duas instâncias de cada agente, e o Jason numera os nomes: \cod{initiator1}, \cod{initiator2}, \cod{participant1} e \cod{participant2}. Cada initiator começa com a crença \cod{tarefas(3)}, que é o $i$. O ambiente recebe ($n$, $i$, $m$) para saber quantas negociações esperar.

Os agentes seguem a divisão do AgentSpeak em crenças, regras e planos. O participant sabe que serviço oferece (\cod{servico(servico\_a)}) e deduz, por regras, o próprio número e a sua estratégia. O initiator começa com o objetivo \cod{!start} e trata as mensagens recebidas com planos.

\subsection*{2.2 Negociações em paralelo}

O initiator abre as $i$ negociações assim:

\begin{lstlisting}
+!start : tarefas(I)
   <- .wait(1000);
      for (.range(K, 1, I)) {
         !!cnp(K)
      }.
\end{lstlisting}

O \cod{.wait(1000)} garante que todos os agentes já tenham sido criados. O ponto principal é o operador \cod{!!}. Com \cod{!cnp(K)}, cada negociação seria um subobjetivo da intenção atual, e o initiator só abriria a segunda depois de terminar a primeira. Com \cod{!!cnp(K)}, cada negociação torna-se uma intenção independente, e as $i$ negociações executam ao mesmo tempo.

Como as negociações executam ao mesmo tempo, elas precisam de um identificador. Cada uma recebe um \cod{Cid = c(Me, K)}, por exemplo \cod{c(initiator1, 3)} para a terceira negociação do initiator1. Toda mensagem e toda crença de uma negociação carregam o seu \cod{Cid}. Sem isso, a resposta de uma negociação poderia ser tomada como resposta de outra.

\subsection*{2.3 Estados de uma negociação}

Cada negociação passa pelos estados abaixo, guardados na crença \cod{estado(Cid, X)}:

\begin{lstlisting}
coletando --> decidir -+-> sem_proposta        (ninguem propos)
                       +-> esperando -+-> concluida          (inform_done do vencedor)
                                      +-> timeout_resultado  (5 s sem entrega)
\end{lstlisting}

O initiator decide assim que todos os participants respondem. Cada proposta ou recusa recebida dispara um plano que verifica se todas já chegaram:

\begin{lstlisting}
todos_responderam(Cid) :-
    cfps(Cid, NP) &
    .count(propose(Cid, _)[source(_)], NA) &
    .count(refuse(Cid)[source(_)], NR) &
    NA + NR >= NP.

+propose(Cid, _) : estado(Cid, coletando) & todos_responderam(Cid)
   <- !decidir(Cid).
\end{lstlisting}

Há um plano equivalente para \cod{refuse}. O prazo de 2000\,ms continua ativo, mas só é usado quando algum participant não responde. Sem ele, o initiator aguardaria indefinidamente.

Os planos que decidem e que encerram uma negociação são marcados como \cod{[atomic]}. Enquanto um deles executa, nenhuma outra intenção do initiator é executada. Assim, uma negociação nunca termina duas vezes: ou chega o \cod{inform\_done} a tempo, ou o prazo expira.

Para trocar o estado, foi usado:

\begin{lstlisting}
+!mudar_estado(Cid, Novo)
   <- -estado(Cid, _);
      +estado(Cid, Novo).
\end{lstlisting}

A forma mais curta seria \cod{-+estado(Cid, Novo)}, mas o \cod{-+} do Jason remove a primeira crença \cod{estado(\_, \_)} que encontrar, e ela pode ser de outra negociação do mesmo initiator. Com várias negociações simultâneas, isso corromperia os estados.

O initiator também verifica a origem de cada mensagem. A entrega só é aceita se vier do vencedor e se a negociação ainda estiver esperando:

\begin{lstlisting}
+inform_done(Cid)[source(W)] : vencedor(Cid, W, Preco) & estado(Cid, esperando)
\end{lstlisting}

O participant, por sua vez, só executa o serviço se tiver feito proposta para aquele initiator naquela negociação, e uma única vez.

\subsection*{2.4 Estratégias de preço}

A estratégia de cada participant é definida pelo seu número, pelo resto da divisão por 3:

\begin{center}
\begin{tabular}{llr}
\toprule
Número mod 3 & Estratégia & Preço \\
\midrule
0 & fixa & sempre 1000 \\
1 & agressiva & sorteado entre 900 e 1050 \\
2 & aleatória & sorteado entre 800 e 1200 \\
\bottomrule
\end{tabular}
\end{center}

Os intervalos se sobrepõem, portanto qualquer estratégia pode vencer. A fixa é previsível. A agressiva quase sempre fica abaixo da fixa, mas nunca desce de 900. A aleatória pode ser a mais barata ou a mais cara de todas. O preço é sorteado novamente a cada pedido. O initiator escolhe o menor preço, e o empate é decidido pelo menor nome. Isso decorre do formato das propostas: o initiator reúne todas em uma lista de \cod{of(Preco, Nome)}, e o \cod{.min} do Jason compara primeiro o preço e depois o nome.

\subsection*{2.5 Ambiente e contagem de mensagens}

O ambiente (\cod{CnpEnv.java}) só aceita uma ação, \cod{registrar}. Os initiators a chamam quando abrem uma negociação e quando ela termina. Cada chamada gera uma linha em \cod{results/eventos.csv}, com o horário em milissegundos, o agente, o evento, o \cod{Cid} e, nas conclusões, o vencedor e o preço. O método é \cod{synchronized}, porque muitos agentes registram ao mesmo tempo. Quando o número de negociações terminadas chega a $n$ $\times$ $i$, o ambiente grava um \cod{resumo.csv} e encerra o processo.

As mensagens são contadas pelo próprio initiator, em cada negociação. No momento da decisão, ele soma os \cod{cfp} que enviou com as propostas e as recusas que chegaram no prazo:

\begin{lstlisting}
      .count(refuse(Cid)[source(_)], NR);
      .length(L, NL);
      ?cfps(Cid, NP);
      T = NP + NL + NR;
      +msgs(Cid, T);
\end{lstlisting}

Depois acrescenta o aceite e as rejeições, uma mensagem para cada proposta recebida, e, se o vencedor entregar, mais uma do \cod{inform\_done}. O total vai no último argumento do \cod{registrar} do estado final, e o ambiente soma os totais de todas as negociações.

Essa contagem considera apenas o que acontece dentro da negociação. Uma proposta que chega depois da decisão, ou um \cod{inform\_done} que chega depois do timeout, não é contabilizada.

\separador

\section*{3. Método dos experimentos}

Foi variado um parâmetro de cada vez, com 3 repetições por cenário:

\begin{center}
\begin{tabular}{lrrr}
\toprule
Bloco & n & m & i \\
\midrule
Variar n & 2, 10, 50, 100, 199 & 10 & 3 \\
Variar m & 50 & 2, 5, 10, 25, 49 & 3 \\
Variar i & 50 & 10 & 1, 3, 5, 9 \\
Estresse & 199 & 49 & 9 \\
\bottomrule
\end{tabular}
\end{center}

O cenário 50/10/3 aparece nos três primeiros blocos e foi executado uma única vez. São 13 cenários e 39 execuções.

As métricas, calculadas a partir do \cod{eventos.csv} de cada execução, são:

\begin{center}
\begin{tabular}{l>{\raggedright\arraybackslash}p{5.2cm}}
\toprule
Métrica & Como é calculada \\
\midrule
Tempo total & da primeira abertura até o último estado final \\
Latência & da abertura até o estado final de cada negociação (média) \\
Vazão & negociações concluídas por segundo de tempo total \\
Taxa de sucesso & concluídas $\div$ ($n$ $\times$ $i$) \\
Mensagens & soma das mensagens de cada negociação, contadas pelo initiator \\
Memória & pico de memória do processo, medido pelo \cod{/usr/bin/time -l} \\
\bottomrule
\end{tabular}
\end{center}

A latência inclui os 100\,ms que o vencedor leva para executar o serviço. O restante é o tempo gasto com as mensagens e com o raciocínio dos agentes.

O tempo total não inclui a inicialização da JVM nem o segundo de espera do \cod{!start}, porque começa a contar na primeira abertura.

\separador

\section*{4. Experimento 1: variação de $n$}

Com $m$ = 10 e $i$ = 3:

\begin{center}
\small
\begin{tabular}{lrrrrrr}
\toprule
n & Tempo total & Desvio & Latência & Vazão & Mensagens & Memória \\
 & (ms) & (ms) & (ms) & (/s) &  & (MB) \\
\midrule
2 & 145 & 3 & 137 & 41,4 & 186 & 94 \\
10 & 205 & 1 & 177 & 146,6 & 930 & 122 \\
50 & 666 & 113 & 572 & 229,3 & 4.650 & 184 \\
100 & 849 & 241 & 749 & 372,4 & 9.300 & 336 \\
199 & 1.919 & 179 & 1.651 & 312,8 & 18.507 & 507 \\
\bottomrule
\end{tabular}
\end{center}

\textbf{Comentário}

Todas as negociações foram concluídas, e o número de mensagens coincidiu com a fórmula em todos os cenários. A latência subiu de 137\,ms com 2 initiators para 1,7\,s com 199.

A vazão cresceu até 100 initiators e caiu com 199, de 372,4 para 312,8 negociações por segundo. A partir desse ponto, mais initiators apenas aumentam o tempo de espera. A memória cresceu cerca de 2,1\,MB por agente.

\separador

\section*{5. Experimento 2: variação de $m$}

Com $n$ = 50 e $i$ = 3:

\begin{center}
\small
\begin{tabular}{lrrrrrr}
\toprule
m & Tempo total & Desvio & Latência & Vazão & Mensagens & Memória \\
 & (ms) & (ms) & (ms) & (/s) &  & (MB) \\
\midrule
2 & 267 & 3 & 223 & 562,5 & 1.050 & 148 \\
5 & 338 & 10 & 277 & 443,6 & 2.400 & 175 \\
10 & 666 & 113 & 572 & 229,3 & 4.650 & 184 \\
25 & 1.010 & 48 & 887 & 148,7 & 11.400 & 312 \\
49 & 1.815 & 312 & 1.690 & 84,4 & 22.200 & 447 \\
\bottomrule
\end{tabular}
\end{center}

\textbf{Comentário}

Cada negociação troca 3$m$ + 1 mensagens, portanto o $m$ é o parâmetro de maior impacto por negociação. De 2 para 49 participants, as mensagens aumentaram 21 vezes, a latência 7,6 vezes, e a vazão caiu de 562,5 para 84,4 negociações por segundo.

\separador

\section*{6. Experimento 3: variação de $i$}

Com $n$ = 50 e $m$ = 10:

\begin{center}
\small
\begin{tabular}{lrrrrrr}
\toprule
i & Tempo total & Desvio & Latência & Vazão & Mensagens & Memória \\
 & (ms) & (ms) & (ms) & (/s) &  & (MB) \\
\midrule
1 & 249 & 9 & 216 & 201,0 & 1.550 & 157 \\
3 & 666 & 113 & 572 & 229,3 & 4.650 & 184 \\
5 & 847 & 5 & 733 & 295,2 & 7.750 & 292 \\
9 & 1.349 & 220 & 1.143 & 339,7 & 13.950 & 425 \\
\bottomrule
\end{tabular}
\end{center}

\textbf{Comentário}

As negociações de cada initiator executaram em paralelo, mas concorrem pelos mesmos participants. Com 9 vezes mais negociações, a vazão aumentou apenas 1,7 vezes, e a latência passou de 216\,ms para 1,1\,s.

Considerando os Experimentos 1 e 3, a latência acompanha o total de negociações ($n$ $\times$ $i$), e não a forma como elas se dividem entre os initiators:

\begin{center}
\begin{tabular}{llr}
\toprule
Negociações & Cenário & Latência \\
 &  & (ms) \\
\midrule
150 & 50 $\times$ 3 & 572 \\
250 & 50 $\times$ 5 & 733 \\
300 & 100 $\times$ 3 & 749 \\
450 & 50 $\times$ 9 & 1.143 \\
597 & 199 $\times$ 3 & 1.651 \\
\bottomrule
\end{tabular}
\end{center}

A Figura~\ref{fig:latencia} mostra a latência nos três experimentos.

\begin{figure}[h]
\centering
\begin{tikzpicture}
\begin{axis}[width=5.2cm, height=4.6cm, xlabel={$n$}, ylabel={latência (ms)},
  ymin=0, ymax=1700, grid=major, grid style={gray!30},
  tick label style={font=\footnotesize}, label style={font=\small}, title={\small $m=10$, $i=3$}]
\addplot[mark=*, thick] coordinates {(2,137) (10,177) (50,572) (100,749) (199,1651)};
\end{axis}
\end{tikzpicture}
\hfill
\begin{tikzpicture}
\begin{axis}[width=5.2cm, height=4.6cm, xlabel={$m$}, ylabel={},
  ymin=0, ymax=1700, grid=major, grid style={gray!30},
  tick label style={font=\footnotesize}, label style={font=\small}, title={\small $n=50$, $i=3$}]
\addplot[mark=*, thick] coordinates {(2,223) (5,277) (10,572) (25,887) (49,1690)};
\end{axis}
\end{tikzpicture}
\hfill
\begin{tikzpicture}
\begin{axis}[width=5.2cm, height=4.6cm, xlabel={$i$}, ylabel={},
  ymin=0, ymax=1700, grid=major, grid style={gray!30},
  tick label style={font=\footnotesize}, label style={font=\small}, title={\small $n=50$, $m=10$}]
\addplot[mark=*, thick] coordinates {(1,216) (3,572) (5,733) (9,1143)};
\end{axis}
\end{tikzpicture}
\caption{Latência média por negociação em cada experimento. Média de 3 repetições.}
\label{fig:latencia}
\end{figure}

\separador

\section*{7. Experimento 4: estresse}

Com $n$ = 199, $m$ = 49 e $i$ = 9, são 248 agentes e 1.791 negociações por execução:

\begin{center}
\small
\begin{tabular}{lrrrrrr}
\toprule
Repetição & Concluídas & Timeout & Taxa de sucesso & Mensagens & Tempo total & Memória \\
 &  &  &  &  & (ms) & (MB) \\
\midrule
1 & 1.644 & 147 & 91,8\% & 250.719 & 24.506 & 1.043 \\
2 & 1.539 & 252 & 85,9\% & 202.860 & 26.423 & 1.255 \\
3 & 1.464 & 327 & 81,7\% & 262.621 & 34.874 & 1.102 \\
\bottomrule
\end{tabular}
\end{center}

A fórmula prevê 265.068 mensagens por execução.

\textbf{Comentário}

Todas as 1.791 negociações terminaram, e de 81,7\% a 91,8\% foram concluídas. As falhas foram todas \cod{timeout\_resultado}. O sistema saturou: as negociações levaram de 10 a 35\,s, e a lentidão fez com que vários prazos de 2 e de 5 segundos expirassem.

O número de mensagens ficou abaixo da fórmula porque, depois de um prazo expirado, as respostas atrasadas não são contabilizadas. Os contratos ficaram distribuídos entre 25 e 35 participants por repetição.

\separador

\section*{8. Avaliação da abordagem de agentes}

Pontos positivos:

\begin{itemize}
  \item Cada papel do CNP corresponde a um tipo de agente, e cada passo do protocolo a uma mensagem. Não foi necessário nenhum coordenador central.
  \item Cada agente armazena apenas o que precisa nas suas crenças. As verificações do protocolo, como aceitar a entrega só se vier do vencedor, ficaram com uma linha cada.
\end{itemize}

Dificuldades encontradas:

\begin{itemize}
  \item A concorrência continuou exigindo cuidado. Foi necessário marcar alguns planos como \cod{[atomic]} e trocar o \cod{-+estado}, como explicado na Seção 2.3.
  \item O Jason agrupa propostas iguais em uma única crença, com várias fontes. Para verificar se todos já haviam respondido, foi necessário contar as fontes (\cod{[source(\_)]}), e não as crenças.
\end{itemize}

\separador

\section*{9. Avaliação das ferramentas}

O Jason mostrou-se adequado para o problema. Os dois agentes têm, juntos, 102 linhas de código, e as ações internas \cod{.send}, \cod{.findall}, \cod{.count}, \cod{.min} e \cod{.wait} atenderam às necessidades do protocolo. O ambiente em Java exigiu apenas dois métodos, \cod{init} e \cod{executeAction}.

\separador

\section*{10. Critério de comparação entre implementações}

Para comparar esta implementação com outra, são propostas duas etapas.

A primeira é a correção: a implementação deve passar nos seis cenários de teste do repositório (\cod{scripts/testes.py}).

A segunda é a eficiência, medida na mesma máquina e nos mesmos cenários:

\begin{center}
\begin{tabular}{lll}
\toprule
Medida & Cenário & Esta implementação \\
\midrule
Latência por negociação & 50, 10, 3 & 572\,ms \\
Latência por negociação & 199, 10, 3 & 1.651\,ms \\
Mensagens por negociação & 50, 10, 3 & 31 (= 3$m$ + 1) \\
Memória por agente & variação de n & 2,1\,MB \\
Taxa de sucesso & 199, 49, 9 & 86,5\% (de 81,7\% a 91,8\%) \\
\bottomrule
\end{tabular}
\end{center}

A latência pode ser comparada diretamente porque o initiator decide assim que todos respondem, sem esperar um prazo fixo. Em uma implementação que sempre espera o prazo, esse prazo deve ser descontado antes da comparação.

\separador

\section*{11. Conclusões}

O trabalho implementou o Contract Net Protocol em Jason e funcionou como esperado. As negociações rodaram em paralelo, cada participant usou a sua estratégia de preço, e o sistema chegou a 597 negociações ao mesmo tempo, com 209 agentes, sem perder nenhuma. Fora do cenário de estresse, todas as negociações foram concluídas e o número de mensagens foi exatamente o que o protocolo prevê.

Nos experimentos, o tempo de cada negociação cresceu com os três parâmetros. O que mais pesou foi o total de negociações ($n$ $\times$ $i$). O $m$ foi o que mais encareceu cada negociação, porque cada participant a mais significa mais mensagens.

No estresse, com 1.791 negociações, o sistema ficou lento, mas todas terminaram, e entre 81,7\% e 91,8\% foram concluídas.

O trabalho tem algumas limitações. Os prazos de 2 e de 5 segundos são fixos e não mudam com a carga, e foi isso que causou os timeouts no estresse. Os participants aceitam todos os contratos que recebem, sem limite. E as medições foram feitas em uma única máquina, com 3 repetições por cenário, o que é pouco para tirar conclusões mais fortes.

\end{document}
